% !PRINTED 404
\weberpage{404}

%<<CONT
haben alle dieselbe Norm; sie sind assoziiert nach $[Q_1]$, aber
nicht nach $[Q]$. Demnach erhalten wir nach (6):

\[ \mathrm{Lim}_{s=1}\sum\frac{s-1}{(Ny_1)^s} = \frac{r}{a\sqrt D}\,\frac{\log\varepsilon}{\lambda}; \]
um daraus die auf die Zahlen (9) bezügliche Summe zu erhalten,
hat man, da $N(\varepsilon_1)=1$, $N(r)=r^2$ ist, mit $\lambda$ zu multiplizieren
und mit $r^2$ zu dividieren. Man erhält so:

\[ \mathrm{Lim}\sum\frac{s-1}{(Ny)^s} = \frac{1}{r\,a\sqrt D}\log\varepsilon. \]
Zieht man dies von der Summe (6) ab und verfährt ebenso
mit allen Primfaktoren $r$ von $Q$, so ergibt sich:

\weq{10}{\mathrm{Lim}\sum\frac{s-1}{(Ny)^s} = \prod\Bigl(1-\frac1r\Bigr)\frac{1}{a\sqrt D}\log\varepsilon, \qquad D>0,}
und noch einfacher bei negativer Diskriminante:

\[ = \prod\Bigl(1-\frac1r\Bigr)\frac{\pi}{a\sqrt{-D}}, \qquad D<0, \]
worin sich die Summe nur auf die Zahlen $y$ erstreckt, die
zu $Q$ teilerfremd sind, und das Produkt $\Pi$ sich auf alle
Primfaktoren von $r$ bezieht. % [sic] Gedruckt ist r; nach dem unmittelbar vorangehenden Satz und der Produktbildung sind die Primfaktoren von Q gemeint.
Nur in den beiden Ausnahmefällen
$D=-4$, $D=-3$ ist die rechte Seite der letzten Formel
noch durch 2 oder durch 3 zu dividieren.

Zerlegt man $y$ in Idealfaktoren

\weq{11}{y = \mathfrak{a}\,\mathfrak{m},}
so durchläuft $\mathfrak{m}$ alle Ideale einer bestimmten Idealklasse nach
der Ordnung $[Q]$, die zu $Q$ teilerfremd sind. Diese Klasse ist,
wenn $\mathfrak{a}$ in die Klasse $A$ gehört,

\[ M = A^{-1}, \]
und unter den Zahlen

\[ N(y) = a\,N(\mathfrak{m}) \]
kommt im Falle der positiven Diskriminante jedes $N(\mathfrak{m})$ nur einmal
, im Falle einer negativen Diskriminante zwei- oder vier- oder
sechsmal vor.

Demnach ergibt sich aus (10):

\weq{12}{\begin{aligned}
\mathrm{Lim}_{s=1}\sum\frac{s-1}{(N\mathfrak{m})^s} &= \prod\Bigl(1-\frac1r\Bigr)\frac{1}{\sqrt D}\log\varepsilon, & D&>0\\
&= \prod\Bigl(1-\frac1r\Bigr)\frac{\pi}{\sqrt{-D}}, & D&<0
\end{aligned}}
%CONT>>
